% Ad Atticum 13.22 --- headnote
% ad-atticum-13-22, 4 July 45 BC, Arpinum

\textit{Cicero to Atticus, written from his Arpinum
estate on 4 July 709 AUC (Perseus dateline
\textit{iv Non.\ Quint.\ a.\ 709 (45)}). Cicero is
still at Arpinum putting the small estates in order
and is about to come back to Tusculanum, where he and
Atticus have arranged to meet. The letter is a
miscellany of literary and financial business: the
\textit{Academica} that Cicero is dedicating to Varro
(the running anxiety of these weeks --- whether Varro
will take the dedication well, and which copies are
fit to be released); the murder of Marcus Marcellus
in the Piraeus, of which Cassius had sent a notice
and Servius Sulpicius Rufus has now sent the
particulars from Athens; the dunning of assigned
debtors; the ongoing Brutus--Claudia trouble; the
embarrassment over Caerellia's having obtained a
draft of the \textit{Academica} from somewhere; and
the proposed purchase of a grove (the planned shrine
to Tullia) which Cicero is starting to argue himself
out of.}

\textit{Three Greek tags, all in the standard Atticus
register: \textit{[Greek: asmenaitata]} ``most
gladly'' (of weaving Atticus into the dialogue),
\textit{[Greek: ta kata meros]} ``the particulars'' (of
Servius' fuller account of Marcellus' death), and
\textit{[Greek: eulogian]} ``plausibility'' (the
intellectual cover the grove proposal has, even if
the location is too remote). The textual crux at the
end of section 4 --- \dag{}\textit{a quis sine te
opprimi militia est}\dag{} --- is preserved as a
dagger; the sense is that Cicero finds it hard going
to handle the co-heirs without Atticus at his elbow.
``Tullius my secretary'' (\textit{Tullium scribam}) is
a freedman scribe of Cicero's, not the orator's son.
The closing complaint that Attica has not sent her
grandfather-in-letters so much as a greeting is one
of the recurrent affectionate teases of the
correspondence.}
